\section{The exact field of definition of the geometric components}
\label{sec:actual-component-field}

The good cover requires a small field over which every original
component is geometrically integral.  The crucial arithmetic hypothesis
is already present in the host: the CM extension $E/E^+$ is unramified
at every finite prime.  We use it without imposing a parity condition
on the rank or replacing the original lattice.

Write $F=E^+$, $K_0=\mathbb Q(\sqrt{-d})\subset E$, and retain the
original trace-selfdual $\mathcal O_E$-lattice $\Lambda$, alternating
form, primitive line and signature $(p,q)$.  Put $n=p+q\geq1$.
The signature is constant above the chosen embedding of $K_0$.
Let $l$ be the original level, $m=\operatorname{lcm}(l.M,l.A)$,
and let $C$ be its original coefficient field.  In particular
$m\geq3$ and $K_0\subset C$.  Denote by $R_m$ the ray class field
of $K_0$ modulo $m\mathcal O_{K_0}$, and put $C'=CR_m$.

\begin{theorem}[Actual component-field supplier]
\label{thm:actual-component-field}
Every geometric irreducible component of the original moduli scheme
is defined and geometrically integral over $C'$.  The exact source
condition $\mathrm{FnAlgClosed}$ therefore holds, including after
adjoining its original generic flag variables.  The field $C'/C$ is
normal, contains the original roots of unity, and satisfies
\[
 \operatorname{fieldSz}(C')\leq\operatorname{fieldSz}(C)
                  +20(\log m+\log d+1),
 \qquad
 \operatorname{fieldSz}(L)=\log[L:\mathbb Q]+\log\operatorname{rd}(L).
\]
The statement includes both definite signatures and all ranks.
For the actual cover, $l.A=1$ and $m=g_0g_1\,15^e$, so its
field-size bound is precisely the prescribed one.
\end{theorem}

\begin{proof}
\paragraph{The integral local calculation.}
Choose the original nonzero anti-invariant $\xi\in E$, for example
$\sqrt{-d}$, and write the alternating form as
\[
                 B(x,y)=\operatorname{Tr}_{E/\mathbb Q}
                                      (\xi h(x,y)).
\]
The adjoint coefficient action determines the nondegenerate Hermitian
form $h$ uniquely.  At a finite place of $F$, trace self-duality gives
\[
 \Lambda=J\Lambda^{\vee,h},\qquad
 J=\xi^{-1}\mathfrak D_{E/\mathbb Q}^{-1}.
                                                        \tag{CF1}
\]
Indeed one may multiply either lattice argument by every local
coefficient integer, so the trace-dual condition is exactly membership
in the indicated codifferent.  This also proves (CF1) at primes over
two; no division by two occurs.

At an unramified quadratic place, the conjugation-stable fractional
ideal $J$ is generated by a scalar $c\in F_v^\times$.
The form $c^{-1}h$ on the actual lattice is unimodular.  Such a form
has an orthogonal integral basis with unit diagonal entries.  To see
this also in residue characteristic two, a nondegenerate Hermitian
form over the quadratic residue field has a vector of nonzero norm:
if all basis norms vanish, a nonzero off-diagonal pairing and the
surjective residue trace give one.  Lift this vector, split off its
unit-norm line integrally, and induct.  The norm on units of an
unramified local extension is onto\cite{component-field-cft}; rescale
each basis vector to make every diagonal entry one.
At a split place the two coefficient factors are paired perfectly,
and dual bases give the corresponding assertion directly.

Two original trace-selfdual lattices of the same signature consequently
give locally isometric rational Hermitian spaces at every finite place.
Their common factor $c$ in (CF1) depends only on $\xi$ and the local
different, so it is the same in both calculations.  They are isometric
at infinity by their prescribed signatures.  The local-to-global
principle for the unitary group now gives a rational isometry in every
rank.  Explicitly, its derived group is the simply connected special
unitary group, and its torus quotient is the restriction of scalars
of the norm-one torus.  The kernel of localization on the first
cohomology of that torus is zero by the quadratic Hasse norm theorem.
The reductive-group Hasse principle in Lemma~8.20 of
\cite{component-field-shimura} applies to the isometry torsor.
For rank one this is exactly the Hasse norm theorem itself.
Thus the rational-isometry condition in the PEL moduli dictionary is
verified, rather than inferred from the name of the moduli problem.

The same calculation gives a second fact that retains the full level:
for every local norm-one unit $v$, there is an integral unitary
automorphism with determinant $v$.  In the orthogonal basis it is
$\operatorname{diag}(v,1,\ldots,1)$; in the split case use the
general-linear factor and its inverse adjoint.  If $v\equiv1\pmod m$,
the automorphism is congruent to the identity on the original lattice
modulo $m$, independently of the chosen integral basis.

\paragraph{Indefinite signatures and the exact reflex map.}
Suppose $p,q>0$.  Let $G$ be the rational unitary similitude group of
the original polarized space, with rational similitude multiplier.
Its derived group is $\operatorname{Res}_{F/\mathbb Q}\mathrm{SU}(h)$,
which is simply connected and noncompact at the real places.
Its commutative quotient and quotient map are
\[
 T=\{(\delta,c)\in\operatorname{Res}_{E/\mathbb Q}\mathbb G_m
                      \times\mathbb G_m:\delta\bar\delta=c^n\},
 \qquad \nu(g)=(\det_E g,c(g)).
                                                        \tag{CF2}
\]
Let $K$ be the actual finite-level stabilizer.  It contains the full
principal congruence subgroup modulo $m$ in the integral similitude
stabilizer.  This containment fixes both complete label maps, hence
also the original selected half-label, and the prescribed pairing
root.  The fixed sign comparison with the generic PEL convention
commutes with these scalar and coefficient-linear transformations.

The Hodge cocharacter has $p$ copies of one $K_0$-embedding and $q$
of its conjugate.  On the torus (CF2) its reflex map, after restriction
to $K_0$, is therefore
\[
        r(s)=(s^p\bar s^q,s\bar s).
                                                        \tag{CF3}
\]
Reversing the reciprocity convention inverts this map and does not
alter the subgroup containment below.  If $p=q$, the reflex field
is smaller and (CF3) factors through $N_{K_0/\mathbb Q}$; restriction
to $K_0$ is still the displayed map.

For $u\in\widehat{\mathcal O}_{K_0}^{\times}$ with
$u\equiv1\pmod m$, scalar multiplication by $u$ on the original
local lattices has multiplier $u\bar u$ and determinant $u^n$.
The remaining determinant factor is
\[
                         v=(\bar u/u)^q.
\]
It is a norm-one unit congruent to one modulo $m$.
The preceding local construction gives an integral unitary matrix
$a_v\equiv I\pmod m$ with this determinant.  Hence
\[
                  ua_v\in K,\qquad \nu(ua_v)=r(u).
                                                        \tag{CF4}
\]
This is an adelic assertion on every original lattice, not the
existence of one global integral orthogonal basis.

The already represented PEL moduli scheme is its canonical Shimura
model with this exact lattice and level.  The preceding rational
isometry check verifies the extra condition in Theorem~8.17 of
\cite{component-field-shimura}; the generic PEL representability
comparison preserves the original primitive line and labels.
The algebraic moduli action is the canonical action by the main
theorem of complex multiplication; this identification is also the
one in Lan's characteristic-zero PEL description
\cite{component-field-pel}.

The connected-component theorem and its reciprocity law give
\[
 \pi_0(\operatorname{Sh}_K)\simeq
       T(\mathbb Q)^\dagger\backslash T(\mathbb A_f)/\nu(K),
 \qquad
 \sigma[t]=[r(s_f)t].                                  \tag{CF5}
\]
Here $\operatorname{art}_{K_0}(s)$ is the restriction of $\sigma$
to $K_0^{\mathrm{ab}}$; see Theorem~5.17 and formula~(64) of
\cite{component-field-shimura}.  The formula for arbitrary $K$
follows by taking a sufficiently small normal subgroup and then its
finite quotient.  That proof does not change the arithmetic level
or coefficient field in the conclusion.
An automorphism fixing $R_m$ has this Artin class represented by
such a $u$.  Equation (CF4) makes its action in (CF5) trivial.
Thus every geometric component is fixed.

\paragraph{The definite endpoints.}
If $q=0$ or $p=0$, the moduli space has dimension zero.
Every corresponding abelian variety is of pure $K_0$-CM type.
It is isogenous to a power of a $K_0$-CM elliptic curve.
Equivalently, a rational orthogonal basis realizes a CM algebra
$E^n$ in its rational endomorphisms; no integral decomposition of
$\Lambda$ is required.  The type on each factor is induced by the
same $K_0$-embedding, so its reflex norm is the central scalar $s$
(or $\bar s$) on the entire rational Tate module.

Let $\sigma$ fix $R_m$ and choose the same unit idele $u$ as above.
The CM theorem for CM algebras gives a quasi-isogeny
$\alpha:A\longrightarrow\sigma A$ satisfying
\[
                         \alpha(ux)=\sigma x
                  \quad\text{for every }x\in V_fA,
                                                        \tag{CF6}
\]
with $\bar u$ in the conjugate-type case
\cite[Theorem~9.10 and Remark~9.11]{component-field-cm}.
Since $u\widehat\Lambda=\widehat\Lambda$, this quasi-isogeny and
its inverse preserve every integral Tate module.  They are actual
inverse isomorphisms: a denominator at a prime would contradict
integrality on that prime's Tate module.  The scalar commutes with
the original $\mathcal O_E$-action, so (CF6) intertwines that action
with its conjugate exactly.

The polarization is also unchanged.  Remark~9.11(c) of the CM
theorem gives the multiplier
\[
 c=\chi_{\mathrm{cyc}}(\sigma)/N_{K_0/\mathbb Q}(u)
                                      \in\mathbb Q^\times.
\]
Both numerator and denominator are finite unit ideles, so
$c\in\mathbb Q^\times\cap\widehat{\mathbb Z}^{\times}
=\{1,-1\}$.  The original and pulled-back forms are
polarizations, so $c>0$ and hence $c=1$.  This retains the
polarization exactly without an Artin-sign convention ambiguity.
Therefore
$\alpha^*(\sigma\lambda)=\lambda$.  The original primitive line
is uniquely recovered, with rigidification, from $\lambda$ and
$w+\bar w=1$, so it too is preserved.  Finally $u\equiv1\pmod m$
makes (CF6) identify every original $M$- and $A$-torsion label.
Every original moduli point is fixed.  This argument includes all
definite lattice classes and does not apply an indefinite
connected-component theorem to a compact derived group.

\paragraph{From components to the exact function fields.}
The original moduli scheme is smooth in characteristic zero.
Its geometric irreducible components are therefore disjoint connected
components.  Their Galois invariance descends each to a geometrically
integral component over $C'$.  Thus $C'$ is algebraically closed in
the function field $\kappa(Z)$ of each component.  The source uses
\[
 \Omega_Z=\operatorname{Frac}
               \bigl(\kappa(Z)[X_\alpha:\alpha\in\mathrm{FlagVar}]\bigr).
\]
A purely transcendental extension introduces no new elements
algebraic over $C'$.  This proves every quantified clause of the
original $\mathrm{FnAlgClosed}$, not merely a statement about
unframed complex points\cite{component-field-source}.

\paragraph{The prescribed field-size bound.}
Put $a=[R_m:K_0]$.  The ray-class exact sequence gives
$a\leq h(K_0)m^2\leq dm^2$.
For the second inequality count reduced binary quadratic forms:
if $A=\lfloor\sqrt{d/3}\rfloor$, their number is at most
$\sum_{a=1}^{A}(2a+1)=A^2+2A\leq3A^2\leq d$.
The conductor--discriminant formula gives
\[
 N_{K_0/\mathbb Q}\mathfrak D_{R_m/K_0}\leq m^{2a},
 \qquad \operatorname{rd}(R_m)\leq\sqrt d\,m.
\]
Consequently
\[
 \operatorname{fieldSz}(R_m)
   \leq\log2+\tfrac32\log d+3\log m
   \leq20(\log m+\log d+1).
\]
Conjugation preserves the ray modulus, so $R_m/\mathbb Q$ is
normal.  Norm compatibility in class field theory gives
$\zeta_m\in R_m$.  Degree and root-discriminant submultiplicativity
for the compositum now give the stated bound for $C'$ and relative
normality over $C$.  The original coefficient embeddings and roots
are retained by the compositum.  No factor depending on $[E:\mathbb Q]$
has entered this field-size estimate.
\end{proof}

The theorem closes the component-field supplier on the original level
cover.  It does not complete its normalized complex theta dictionary,
short-section bounds, determinant-ideal comparison or global analytic
norms.  In particular it places no bound on the growing geometric
dimension or on integral eigenlattice defects.
